\documentclass[11pt]{amsart}
\usepackage[margin=1.1in]{geometry}
\usepackage{amsmath,amssymb,amsthm,mathtools,enumitem}
\usepackage[dvipsnames]{xcolor}
\usepackage[colorlinks=true,linkcolor=blue!50!black,citecolor=blue!50!black]{hyperref}

\newtheorem{theorem}{Theorem}[section]
\newtheorem{proposition}[theorem]{Proposition}
\newtheorem{lemma}[theorem]{Lemma}
\newtheorem{corollary}[theorem]{Corollary}
\theoremstyle{definition}
\newtheorem{definition}[theorem]{Definition}
\theoremstyle{remark}
\newtheorem{remark}[theorem]{Remark}

\newcommand{\R}{\mathbb{R}}

\title{The chromatic number of the plane is seven}
\author{Ansh Mishra}
\author{Aryan Senthilkumar}
\date{\today}

\begin{document}
\begin{abstract}
We prove that every coloring of the Euclidean plane with six colors contains two
points at distance one with the same color. Consequently the chromatic number of
the plane equals seven.
\end{abstract}
\maketitle

\section{Introduction}
% TODO

\section{Weak measurable colorings and circle palettes}
% TODO: recall transfer theorem and palette machinery

\section{Triple junctions}
% TODO

\section{Points of higher chromaticity}
% TODO

\section{Pairs of triple junctions}
% TODO

\section{The global argument}
% TODO

\bibliographystyle{amsplain}
\bibliography{refs}
\end{document}
